\section{The divisor class group}\label{sec:cl}
% ===============================================================

\subsection{The degree descends to classes}

\begin{smjproposition}\label{prop:descend}
The degree induces a well-defined group homomorphism
$\overline{\deg} : \Cl(\PP) \to \ZZ$, $[D] \mapsto \deg(D)$.
\end{smjproposition}

\begin{proof}
If $[D] = [E]$ then $D - E \in \Prin(\PP) = \ker(\deg)$ by
Theorem~\ref{thm:main}, so $\deg(D) - \deg(E) = \deg(D-E) = 0$. Hence
$\deg(D)$ depends only on $[D]$. It is a homomorphism because $\deg$ is one and
the quotient map is one.
\end{proof}

\begin{smjproposition}\label{prop:surj}
$\overline{\deg}$ is surjective.
\end{smjproposition}

\begin{proof}
Given $n \in \ZZ$, pick any $p \in \PP$; then $\overline{\deg}([n(p)]) =
\deg(n(p)) = n$.
\end{proof}

\begin{smjproposition}\label{prop:inj}
$\overline{\deg}$ is injective.
\end{smjproposition}

\begin{proof}
A homomorphism of groups is injective if and only if its kernel is trivial, so it
suffices to show $\overline{\deg}([D]) = 0$ forces $[D] = 0$. If
$\overline{\deg}([D]) = 0$ then $\deg(D) = 0$, so $D \in \ker(\deg) =
\Prin(\PP)$ by Theorem~\ref{thm:main}, and hence $[D] = 0$ in the quotient.
\end{proof}

\begin{theorem}\label{thm:cl}
$\overline{\deg} : \Cl(\PP) \to \ZZ$ is an isomorphism; in particular
$\Cl(\PP) \cong \ZZ$.
\end{theorem}

\begin{proof}
By Propositions~\ref{prop:descend}, \ref{prop:surj} and \ref{prop:inj} it is a
well-defined bijective group homomorphism.
\end{proof}

\subsection{The same computation, packaged}

It is worth seeing that the three verifications above are exactly the content of
one standard theorem, since this both shortens the argument and explains what
made it work.

\begin{smjproposition}
$\Cl(\PP) \cong \ZZ$.
\end{smjproposition}

\begin{proof}
By Theorem~\ref{thm:main}, $\Prin(\PP) = \ker(\deg)$, so
\[
\Cl(\PP) = \frac{\Div(\PP)}{\Prin(\PP)} = \frac{\Div(\PP)}{\ker(\deg)}
\;\cong\; \im(\deg)
\]
by the First Isomorphism Theorem. Finally $\deg(n(p)) = n$ for any fixed $p$, so
$\im(\deg) = \ZZ$.
\end{proof}

The point is that well-definedness, injectivity and surjectivity in
Section~5.1 are precisely the three things the First Isomorphism Theorem checks
in one step, once one knows $\Prin(\PP) = \ker(\deg)$. All the mathematical
content sits in Theorem~\ref{thm:main}.

\subsection{Normal form and interpretation}

\begin{smjproposition}\label{prop:normalform}
Every $D \in \Div(\PP)$ satisfies $D \sim (\deg D)(\infty)$, and the integer
$\deg D$ is the unique one with this property.
\end{smjproposition}

\begin{proof}
The divisor $D - (\deg D)(\infty)$ has degree $\deg(D) - \deg(D) = 0$, hence lies
in $\ker(\deg) = \Prin(\PP)$ by Theorem~\ref{thm:main}; that is,
$D \sim (\deg D)(\infty)$. If also $D \sim n(\infty)$ then
$n = \deg(n(\infty)) = \deg(D)$ by Proposition~\ref{prop:descend}.
\end{proof}

\begin{remark}
Note that Proposition~\ref{prop:normalform} is a consequence of
Theorem~\ref{thm:main} and cannot be used in its proof.
\end{remark}

So on $\PP$ a divisor class carries exactly one piece of information: an integer.
Unwinding the definitions, this says something concrete about functions. A
rational function is determined by its zeros and poles up to a constant factor
(Proposition~\ref{prop:kernel}); a prescribed configuration of zeros and poles is
realised by a function exactly when the multiplicities balance
(Theorem~\ref{thm:main}); and what survives of a configuration after allowing
oneself to modify it by any function at all is precisely its total multiplicity.

% ===============================================================
